\section{Discussion}\label{sec:discussion}

\subsection{What the evidence supports}

Source outcomes help when source and target tasks organize safety around
similar policy shapes.  That is the central empirical claim, and the adverse
regimes make it testable rather than automatic.  Source scoring improves
coverage on the full randomized mixture, but the regime table and Energy
study show concrete cases in which it loses.

This interpretation differs from claiming a generic high-dimensional
optimizer.  The raw vector may contain $10^4$ coordinates, but the decision is
an ordered profile and the library has finite structural complexity.  A
controller may still require a fine lookup table, so profile-grid refinement
is operationally relevant.  It is not equivalent to adding $10^4$ unrelated
decision variables.  The proper benefit is invariance to representation
resolution under shared profile structure.

\subsection{What Energy teaches us}

The Energy reversal shows that a generic structural design can beat source
scoring when rankings from held-out regions do not transfer.  Because V3 was
fixed before its outcomes, changing the score now would turn a test case into
development data.  We keep the result as a negative control and as a practical
boundary on when the method should be used.

The practical implication is a diagnostic choice.  If historical tasks and a
new target share an ordered policy semantics and source rank stability can be
validated on genuinely held-out tasks, source scoring is reasonable.  If the
target is temporally or geographically shifted and generic low-frequency
structure already has a credible physical basis, the generic maximin design is
the safer default.  The method should abstain from a transfer claim when these
conditions cannot be assessed.

The current method begins after this assessment: a source archive has already
been supplied, and the algorithm ranks profiles inside it.  It does not
retrieve related tasks from an unrestricted historical database.  Regime-matched
synthetic sources therefore identify the value of scoring a relevant archive,
whereas Energy shows the cost of archive mismatch.  Source retrieval and its
own data-acquisition cost remain separate research and operational decisions.

\subsection{When the source archive pays for itself}

Ten target calls do not beat 394 target calls for a single target.  The source
archive becomes economically plausible only when it is reused.  Let
$C_0=N_0+V_0$ be the expected target-only calls per deployment and
$C_A=N_A+V_A$ the source-design calls excluding the archive.  The call-count
break-even number of deployments is
\begin{equation}\label{eq:break-even}
 M_{\rm break}=\left\lceil\frac{S}{C_0-C_A}\right\rceil,
 \qquad C_0>C_A.
\end{equation}
In the randomized primary experiment, mean-call estimates give thresholds of
about 7, 9, and 11 targets relative to generic DCT, natural structure, and raw
Sobol.  This arithmetic does not include differences in simulator setup,
wall-clock parallelism, or policy quality, and Energy shows that paying the
archive can still be inferior.  It is a planning threshold, not a dominance
theorem.

Table \ref{tab:outcome-adjusted-cost} also reports calls per certified
deployment, but that ratio is only a descriptive planning aid.  A complete
operational comparison would attach decision-maker costs to unsafe deployment,
abstention, and objective quality.  We report those components separately
rather than choose their relative value on the reader's behalf.

\subsection{Search and verification answer different questions}

Independent verification frequently costs an order of magnitude more than
target search because it answers a different question.  Search asks whether a
promising policy has been found.  Verification asks whether a fixed policy can
be deployed with a specified false-certification bound.  A target-only
optimizer making the same safety claim must also pay for fresh samples.

The all-success rule is valid but weak near the threshold, which explains why
feasible coverage exceeds certified deployment.  Future studies should fix a
more efficient exact or sequential test in advance and use block-aware
verification when temporal replications are not iid.

\subsection{Benchmark fit and remaining threats}

Fixing the generator before confirmatory runs reduces development overfitting
but cannot erase it: the task law, library, and stress axes were created within
the project.  Ordered semantics are a strong prior.  Energy covers only five
regions, and the native transfer study has only three historical families.
The archive is assumed to have been supplied; source retrieval is unsolved.
Three replications cannot order near ties, source aggregation can hide task
disagreement, and a finite library can miss the safe region entirely.  Regret
uses a finite evaluation library rather than the unknown continuous optimum,
and one functional-SCBO failure remains in the denominator.  The theory is
conditional on alignment, independence within each regime, and the declared
replication law.  Finally, the same target seeds appear at all three grid
resolutions, so 480-cell totals are descriptive rather than 480 independent
tasks.
